\documentclass[12pt, a4paper]{article}

\usepackage[utf8]{inputenc}
\usepackage[spanish]{babel}
\usepackage{geometry}
\usepackage{titlesec}
\usepackage{enumitem}
\usepackage{xcolor}
\usepackage{listings}
\usepackage{graphicx}
\usepackage{float}
\usepackage{hyperref}
\usepackage{fancyhdr}

\geometry{margin=2.5cm}

% Colores
\definecolor{azulprincipal}{RGB}{30, 90, 160}
\definecolor{verdeexito}{RGB}{76, 175, 80}
\definecolor{rojopeligro}{RGB}{244, 67, 54}
\definecolor{codigofondo}{RGB}{240, 244, 250}

% Estilo de secciones
\titleformat{\section}{\Large\bfseries\color{azulprincipal}}{}{0em}{}[\titlerule]
\titleformat{\subsection}{\large\bfseries\color{azulprincipal}}{}{0em}{}

% Header y Footer
\pagestyle{fancy}
\fancyhf{}
\rhead{\textit{Space Invaders}}
\lhead{\textit{Grupo Cachopin}}
\cfoot{\thepage}

% --------------------------------------------------
\begin{document}

% PORTADA
\begin{titlepage}
    \centering
    \vspace*{2cm}
    
    {\LARGE \textbf{SPACE INVADERS}}\\[0.5em]
    {\Large \textit{Defensa Planetaria contra la Invasión Alienígena}}\\[1.5em]
    
    \vspace{1cm}
    
    {\large \textbf{Proyecto de Programación}}\\[0.3em]
    {\large Programación Orientada a Objetos — Java}\\[1em]
    
    \vspace{2cm}
    
    {\textbf{Grupo: Cachopin}}\\[0.3em]
    {\large \today}\\[3em]
    
    \vspace{3cm}
    
    {\large \textit{Una batalla épica en pixeles donde la humanidad depende de tu destreza.}}
    
    \vfill
    
    \includegraphics[width=3cm]{SpaceInvaders.png}
    
\end{titlepage}

\newpage
\tableofcontents
\newpage

% ==================================================
\section{Introducción}

\subsection{Presentación del Proyecto}

\textbf{Space Invaders} es un proyecto educativo de programación orientada a objetos que implementa el videojuego clásico del mismo nombre. El objetivo es defender nuestro planeta de una invasión alienígena mediante una nave espacial controlada por el jugador, que debe eliminar a todos los enemigos antes de que estos logren llegar al planeta.

Este proyecto no es solo un simple juego, sino una \textbf{herramienta pedagógica} diseñada para enseñar e implementar conceptos fundamentales de arquitectura de software:

\begin{itemize}[leftmargin=2em]
    \item Patrón arquitectónico \textbf{Modelo-Vista-Controlador (MVC)}
    \item Patrón de diseño \textbf{Observer} para sincronización entre componentes
    \item Patrones de creación y comportamiento (\textbf{Factory}, \textbf{Strategy}, \textbf{Composite})
    \item Programación estructurada con \textbf{SCRUM} y metodología ágil
    \item Gestión de versiones con \textbf{Git} y colaboración en equipo
\end{itemize}

El proyecto está dividido en \textbf{tres sprints} que cubren gradualmente los requisitos, desde funcionalidad básica hasta extensiones avanzadas, permitiendo un desarrollo iterativo y bien planificado.

% ==================================================
\section{Descripción del Juego}

\subsection{Mecánica Básica}

Space Invaders es un juego de defensa donde:

\begin{description}[leftmargin=2em, style=nextline]
    \item[\textbf{Objetivo Principal}]
        Controlar una nave espacial para eliminar todos los enemigos que descienden desde la parte superior de la pantalla, protegiéndola del impacto directo con los enemigos y evitando que estos lleguen al fondo de la matriz.
    
    \item[\textbf{Área de Juego}]
        Una matriz de \textbf{100 × 60 píxeles} que representa el espacio donde ocurre la batalla. El jugador comienza posicionado en el píxel (50, 55), cerca del centro inferior de la pantalla.
    
    \item[\textbf{Enemigos}]
        Se generan entre 4 y 8 enemigos de forma aleatoria en la fila superior (posiciones y entre 0-4), sin que se superpongan entre ellos. Descienden hacia el jugador a razón de una posición cada 200 milisegundos.
    
    \item[\textbf{Movimiento del Jugador}]
        La nave del jugador se mueve libremente mediante las \textbf{flechas direccionales} (arriba, abajo, izquierda, derecha), pero no puede salir de los límites de la matriz.
    
    \item[\textbf{Sistema de Disparos}]
        El jugador dispara mediante la barra \textbf{espaciadora}. Los disparos suben hacia los enemigos a razón de 1 píxel cada 50 milisegundos. Cada disparo sigue su trayectoria independientemente de los otros.
\end{description}

\subsection{Naves Disponibles}

El jugador puede elegir entre tres naves diferentes, cada una con características únicas:

\begin{description}[leftmargin=2em, style=nextline]
    \item[\textbf{Nave Green}]
        Nave de complejidad básica. Tiene acceso a dos tipos de disparo: \texttt{Pixel} (ilimitado) y \texttt{Flecha} (30 disparos). Ideal para principiantes.
    
    \item[\textbf{Nave Blue}]
        Nave de complejidad media. Tiene acceso a dos tipos de disparo: \texttt{Pixel} (ilimitado) y \texttt{Rombo} (20 disparos). Ofrece mayor potencia defensiva.
    
    \item[\textbf{Nave Red}]
        Nave de máxima complejidad. Tiene acceso a los tres tipos de disparo: \texttt{Pixel} (ilimitado), \texttt{Flecha} (30 disparos) y \texttt{Rombo} (20 disparos). Para jugadores avanzados.
\end{description}

\subsection{Tipos de Disparo}

El sistema de disparos utiliza el patrón \textit{Strategy} para permitir múltiples tipos intercambiables:

\begin{description}[leftmargin=2em, style=nextline]
    \item[\textbf{Pixel (Básico)}]
        Un único píxel. Disparos ilimitados. Es el arma estándar siempre disponible. Perfecto para economía de munición.
    
    \item[\textbf{Flecha (Intermedio)}]
        Forma de V invertida compuesta por 3 píxeles. Disponible: 30 disparos. Proporciona mejor cobertura y mayor poder destructivo.
    
    \item[\textbf{Rombo (Avanzado)}]
        Forma de rombo compuesta por 5 píxeles. Disponible: 20 disparos. Es el arma más poderosa pero con munición limitada.
\end{description}

El jugador puede cambiar entre tipos de disparo en cualquier momento. El sistema se encarga automáticamente de elegir solo aquellos tipos que tengan munición disponible.

% ==================================================
\section{Reglas del Juego}

\subsection{Reglas de Movimiento}

\begin{enumerate}[leftmargin=2em]
    \item La nave del jugador se mueve mediante las \textbf{flechas direccionales}.
    \item La nave \textbf{no puede salir de los límites} de la matriz (0-99 en X, 0-59 en Y).
    \item Los enemigos \textbf{descienden automáticamente} cada 200 milisegundos.
    \item Los disparos \textbf{suben automáticamente} cada 50 milisegundos.
    \item El movimiento es \textbf{suave y continuo}, sin saltos o comportamientos abruptos.
\end{enumerate}

\subsection{Reglas de Combate}

\begin{enumerate}[leftmargin=2em]
    \item El jugador dispara mediante la \textbf{barra espaciadora}.
    \item Se pueden disparar \textbf{múltiples proyectiles simultáneamente}, cada uno con su trayectoria independiente.
    \item Puede cambiar el tipo de disparo mediante un \textbf{botón del teclado}.
    \item Las restricciones de munición son:
        \begin{itemize}[leftmargin=2em]
            \item \texttt{Pixel}: ilimitado
            \item \texttt{Flecha}: 30 disparos máximo
            \item \texttt{Rombo}: 20 disparos máximo
        \end{itemize}
    \item Cuando se agota la munición de un tipo, \texttt{no se puede disparar} con ese tipo hasta reinicar la partida.
\end{enumerate}

\subsection{Reglas de Colisión}

\begin{enumerate}[leftmargin=2em]
    \item Si un \textbf{disparo toca a un enemigo}, ambos \textbf{se destruyen} simultáneamente.
    \item Si un \textbf{enemigo toca la nave}, el \textbf{juego termina en derrota}.
    \item Si un \textbf{enemigo alcanza el fondo} (y = 60), el juego termina en derrota y \textit{la invasión terrestre ha comenzado}.
\end{enumerate}

% ==================================================
\section{Condiciones de Victoria y Derrota}

\subsection{Condiciones de Derrota (Game Over)}

El jugador \textbf{pierde la partida} en cualquiera de estos casos:

\begin{itemize}[leftmargin=2em]
    \item Un enemigo entra en \textbf{contacto directo} con la nave del jugador.
    \item Un enemigo \textbf{alcanza la parte inferior} de la matriz sin ser eliminado.
\end{itemize}

En ambos casos, \textit{el enemigo ha llegado a la tierra y la invasión terrestre habrá comenzado}. La humanidad ha fallado.

\subsection{Condición de Victoria}

El jugador \textbf{gana la partida} si:

\begin{itemize}[leftmargin=2em]
    \item \textbf{Elimina todos los enemigos} presentes en pantalla.
\end{itemize}

Cuando esto ocurre, \textit{la humanidad se salva}. El jugador ha defendido exitosamente el planeta.

% ==================================================
\section{Objetivos del Proyecto}

\subsection{Objetivos Técnicos}

\subsubsection{Arquitectura}

\begin{itemize}[leftmargin=2em]
    \item Implementar el patrón \textbf{MVC (Modelo-Vista-Controlador)} de forma perfecta en todos los sprints.
    \item Integrar el patrón \textbf{Observer} para sincronización automática entre modelo y vistas.
    \item Aplicar patrones de diseño: \textbf{Factory}, \textbf{Strategy} y \textbf{Composite}.
    \item Mantener \textbf{bajo acoplamiento} entre componentes y \textbf{alta cohesión} dentro de cada módulo.
    \item Garantizar que \textbf{el modelo no conoce la vista} y viceversa.
\end{itemize}

\subsubsection{Funcionalidad Core}

\begin{itemize}[leftmargin=2em]
    \item Desarrollar un \textbf{motor de juego} con bucle principal (game loop) a 50ms.
    \item Implementar \textbf{sistema de movimiento} para nave, enemigos y disparos.
    \item Crear \textbf{detección de colisiones} eficiente y precisa.
    \item Gestionar \textbf{múltiples enemigos} en pantalla simultáneamente.
    \item Implementar \textbf{diferentes tipos de naves} con características distintas.
    \item Desarrollar \textbf{sistema de disparos intercambiable} con múltiples estrategias.
\end{itemize}

\subsubsection{Interfaz de Usuario}

\begin{itemize}[leftmargin=2em]
    \item Pantalla de inicio clara y funcional.
    \item Pantalla de juego con visualización en cuadrícula de 100×60 píxeles.
    \item Feedback visual claro (colores diferenciados para nave, enemigos, disparos).
    \item Mensajes de victoria y derrota.
    \item Controles intuitivos mediante teclado.
\end{itemize}

\subsection{Objetivos Educativos}

\begin{itemize}[leftmargin=2em]
    \item Comprender y aplicar correctamente el patrón \textbf{MVC} en un proyecto real.
    \item Aprender cómo el patrón \textbf{Observer} permite comunicación desacoplada.
    \item Dominar patrones de creación (\textbf{Factory}) y comportamiento (\textbf{Strategy}, \textbf{Composite}).
    \item Desarrollar habilidades de \textbf{programación orientada a objetos}.
    \item Practicar \textbf{metodología SCRUM} y desarrollo iterativo.
    \item Aprender a usar \textbf{Git} y control de versiones en equipo.
    \item Mejorar capacidades de \textbf{diseño de software} y arquitectura.
\end{itemize}

\subsection{Objetivos de Proyecto}

\begin{itemize}[leftmargin=2em]
    \item Completar los \textbf{tres sprints} dentro de los plazos establecidos.
    \item Entregar un \textbf{producto funcionalmente completo} que cumpla todos los requisitos.
    \item Generar \textbf{documentación de calidad} explicando decisiones de diseño.
    \item Mantener \textbf{colaboración efectiva} del equipo a través de Git.
    \item Presentar el proyecto de forma profesional al cliente.
    \item Opcionalmente, implementar extensiones (\textbf{HU11}) para obtener calificación máxima.
\end{itemize}

% ==================================================
\section{Estructura del Proyecto}

\subsection{Sprints}

El proyecto se desarrolla en \textbf{tres sprints} de evolución incremental:

\begin{description}[leftmargin=2em, style=nextline]
    \item[\textbf{Sprint 1 (Semana 9)}]
        \textit{Funcionalidad Básica}
        \begin{enumerate}[leftmargin=1.5em]
            \item Implementar la pantalla de inicio
            \item Inicializar la matriz del espacio
            \item Implementar nave de un píxel y disparo de un píxel
            \item Implementar enemigo de un píxel
            \item Aplicar perfectamente el patrón MVC
            \item \textbf{Entregable:} Juego funcional básico con MVC
        \end{enumerate}
    
    \item[\textbf{Sprint 2 (Semana 12)}]
        \textit{Patrones de Diseño y Flexibilidad}
        \begin{enumerate}[leftmargin=1.5em]
            \item Ampliar selección de naves (Green, Blue, Red)
            \item Implementar naves y enemigos multipíxel
            \item Implementar tipos de disparo intercambiables
            \item Aplicar patrones de diseño: Factory, Strategy, Composite
            \item \textbf{Entregable:} Sistema de naves y disparos flexible y escalable
        \end{enumerate}
    
    \item[\textbf{Sprint 3 (Semana 15)}]
        \textit{Extensiones y Presentación Final}
        \begin{enumerate}[leftmargin=1.5em]
            \item Agregar características Java 8
            \item Implementar extensiones (HU11) — opcional pero recomendado
            \item Generar documentación final completa
            \item Preparar demo day y presentación profesional
            \item \textbf{Entregable:} Producto final pulido con documentación
        \end{enumerate}
\end{description}

\subsection{Entregas Principales}

\begin{itemize}[leftmargin=2em]
    \item \textbf{Sprint 1:} Código funcional, diagrama UML, informe MVC, repositorio Git
    \item \textbf{Sprint 2:} Código funcional, diagrama UML, informe de patrones, repositorio Git
    \item \textbf{Sprint 3:} Documentación final, demo, diagrama UML final, repositorio Git completo
\end{itemize}

% ==================================================
\section{Requisitos de Evaluación}

\subsection{Requisitos Mínimos (Nota Máxima 8.5/10)}

Para aprobar el proyecto con nota máxima de 8.5 puntos:

\begin{itemize}[leftmargin=2em]
    \item \textbf{Trabajo Presencial:} Todo el grupo debe participar de forma presencial y colaborativa.
    
    \item \textbf{Control de Versiones:} Usar Git desde el inicio y demostrar que todos los participantes aportan mediante commits.
    
    \item \textbf{Entrega Funcional:} Entregar un producto que cumpla \textbf{todos los requisitos hasta HU10}.
    
    \item \textbf{Patrón MVC:} El patrón MVC debe estar \textbf{perfectamente implementado en todos los sprints}. Si no está bien implementado en el Sprint 1, \textcolor{rojopeligro}{\textbf{el proyecto se suspenderá}}.
\end{itemize}

\subsection{Extras (Nota Máxima 10/10)}

Para obtener la nota máxima de 10 puntos:

\begin{itemize}[leftmargin=2em]
    \item \textbf{HU11:} Completar satisfactoriamente la tarea HU11 (expansión del juego).
    
    Ejemplos de expansiones posibles:
    \begin{itemize}[leftmargin=2em]
        \item Sistema de puntuaciones y ranking
        \item Enemigos "inteligentes" con disparos propios
        \item Múltiples niveles de dificultad
        \item Jefe final (Boss)
        \item Integración de Modelos de Lenguaje (LLMs)
        \item Mejoras visuales y de sonido
        \item Otras características creativas
    \end{itemize}
\end{itemize}

% ==================================================
\section{Consideraciones Importantes}

\subsection{Trabajo en Equipo}

\begin{itemize}[leftmargin=2em]
    \item El trabajo en grupo es \textbf{presencial y esencial}. La ausencia de algún miembro será penalizada.
    
    \item El equipo es responsable de su gestión. Si un miembro no participa, el equipo puede expulsarlo notificando al profesor.
    
    \item La calificación de proyectos con \textbf{participación inequitativa se multiplicará por 0.6}.
    
    \item Todos los integrantes deben ser capaces de responder cualquier pregunta del cliente sobre el proyecto.
\end{itemize}

\subsection{Git y Commits}

\begin{itemize}[leftmargin=2em]
    \item El proyecto \textbf{debe estar en GitHub} desde el inicio.
    \item Cada alumno justifica su aportación mediante \textbf{commits con su nombre}.
    \item Los commits deben ser \textbf{significativos} y \textbf{descriptivos}.
    \item Historial de commits será verificado en las reuniones con el cliente.
\end{itemize}

\subsection{Reuniones con el Cliente}

Al final de cada sprint:

\begin{itemize}[leftmargin=2em]
    \item \textbf{Duración:} 15 minutos de reunión presencial.
    \item \textbf{Participación:} Todos los miembros del equipo.
    \item \textbf{Contenido:} Demostración del producto, revisión del código, presentación del diseño.
    \item \textbf{Retroalimentación:} Se anotan cambios y mejoras para el siguiente sprint.
    \item \textbf{Entrega:} Realizada en eGela justo antes de la reunión.
\end{itemize}

% ==================================================
\section{Conclusión}

Space Invaders es más que un videojuego: es una \textbf{experiencia de aprendizaje} que integra arquitectura de software, patrones de diseño, metodología ágil y trabajo en equipo. 

Los objetivos van más allá del simple funcionamiento del juego. Buscan que los estudiantes:

\begin{itemize}[leftmargin=2em]
    \item \textbf{Comprendan} conceptos fundamentales de ingeniería de software
    \item \textbf{Apliquen} patrones reconocidos en la industria
    \item \textbf{Colaboren} efectivamente en equipos
    \item \textbf{Documenten} su trabajo de forma profesional
    \item \textbf{Iterem} continuamente sobre el diseño y la implementación
\end{itemize}

Con dedicación, trabajo en equipo y atención al detalle, este proyecto puede convertirse en un \textbf{excelente ejemplo} de desarrollo de software bien estructurado y profesional.

\vspace{2em}

\begin{center}
    \textcolor{verdeexito}{\Large \textbf{¡Que comience la defensa planetaria!}}
\end{center}

\end{document}
